\begin{table}[t]
\centering
\small
\caption{Comparaison avec le benchmark officiel MedMNIST v2 sur le split de test de PneumoniaMNIST (624 images). Les
valeurs publiées sont reprises de Yang \emph{et al.}~\cite{yang2023medmnist}~; le nombre entre parenthèses est la
résolution d'entrée. Nos CNN~: moyenne$\pm$écart-type sur 3 seeds~; nos modèles classiques sont déterministes
(une seule exécution). Acc.~: seuil par défaut de 0,5, comme dans le benchmark~; \emph{recal.}~: seuil re-choisi sur
la validation (hors protocole du benchmark).}
\label{tab:benchmark}
\begin{adjustbox}{max width=\textwidth}
\begin{tabular}{llccc}
\toprule
Source & Modèle & ROC-AUC & Acc. & Acc. (recal.) \\
\midrule
MedMNIST v2~\cite{yang2023medmnist} & ResNet-18 (28) & 0,944 & 0,854 & n.d. \\
 & ResNet-18 (224) & 0,956 & 0,864 & n.d. \\
 & ResNet-50 (28) & 0,948 & 0,854 & n.d. \\
 & ResNet-50 (224) & 0,962 & 0,884 & n.d. \\
 & auto-sklearn & 0,942 & 0,855 & n.d. \\
 & AutoKeras & 0,947 & 0,878 & n.d. \\
 & Google AutoML Vision & \textbf{0,991} & \textbf{0,946} & n.d. \\
\midrule
Ce travail (28) & LogReg-SGA & 0,927 & 0,851 & n.d. \\
 & SVM linéaire & 0,925 & 0,864 & n.d. \\
 & SVM RBF & 0,947 & 0,857 & 0,872 \\
 & CNN softmax & 0,946$\pm$0,001 & 0,779$\pm$0,067 & 0,869$\pm$0,010 \\
 & CNN EDL & 0,940$\pm$0,014 & 0,761$\pm$0,084 & 0,860$\pm$0,011 \\
\bottomrule
\end{tabular}
\end{adjustbox}
\end{table}
